\begin{question}{26}{
    In het kader van logistieke planning en munitiebeheer is een artillerie-eenheid belast met het optimaliseren van vuursteun bij gevechtsoperaties.
    Om deze optimalisaties te kunnen onderbouwen zijn op basis van een steekproef van elf gevechtsoperaties data verzameld voor het aantal verbruikte munitiekisten (onafhankelijke variabele $X$) 
    en het percentage succesvol geneutraliseerde doelen (afhankelijke variabele $Y$). 
    Deze data zijn te vinden in onderstaande tabel
    {\par\bigskip
        \begin{center}
            \begin{tabular}{c|ccccccccccc}
                \toprule
                    \textbf{$X$} & $14.2$ & $27.8$ & $33.5$ & $46.1$ & $52.4$ & $61.9$ & $74.3$ & $85.6$ & $93.2$ & $108.7$ & $119.4$ \\
                \midrule
                    \textbf{$Y$} & $43.5$ & $26.2$ & $41.1$ & $64.4$ & $28.2$ & $45.6$ & $66.8$ & $62.3$ & $89.5$ & $71.2$ & $94.6$ \\
                \bottomrule
            \end{tabular}
        \end{center}\par\bigskip
    }
}

    % \subquestion{2}{
    %     Als we een regressie-analyse willen uitvoeren, welke variabele zou dan de afhankelijke variabele $Y$ zijn en welke variabele zou de onafhankelijke variabele $X$ zijn?
    % }
    % \solution{
    %     Bij een regressie-analyse is de meest logische keuze om het aantal verbruikte munitiekisten als onafhankelijke variabele $X$ te kiezen,
    %     en het percentage succesvolle geneutraliseerde doelen als afhankelijke variabele $Y$. \rubric{1}
    %     s
    %     Een hoger munitieverbruik zou een hoger percentage succesvol geneutraliseerde doelen kunnen verklaren, andersom is dat minder evident.\rubric{1}
    % }

    \subquestion{4}{
        Teken het spreidingsdiagram (scatter plot) op basis van je antwoord op vraag (a).
        Geef hierbij ook duidelijk aan welke as welke variabele beschrijft.
    }
    \solution{
        De $x$-as duidt de onafhankelijke variabele $X$ uit, oftewel het aantal verbruikte munitiekisten. \rubric{1}

        De $y$-as duidt de afhankelijke variabele $Y$ uit, oftewel het percentage succesvol geneutraliseerde doelen.\rubric{1}
        \begin{center}
            \resizebox{\textwidth}{!}{
                \includegraphics{exam20260724-scatterplot.png} \rubric{2}
            }
        \end{center}
    }

    \subquestion{8}{
        Bereken Pearson's correlatiecoëfficiënt $r(x,y)$ op basis van de data in de bovenstaande tabel.
        Wat zegt $r(x, y)$ over de relatie tussen het cargogewicht en het kerosineverbruik?
    }
    \solution{
        We beginnen met het uitrekenen van Pearson's correlatiecoëfficiënt 
        \begin{align*}
            r(x, y)  &= \frac{ \overline{x \cdot y} - \overline{x} \cdot \overline{y} }{ \sqrt{ (\overline{x^2} - \overline{x}^2 ) \cdot (\overline{y^2} - \overline{y}^2) }}
        \end{align*}
        Hiervoor hebben we dus de waardes van $\overline{x}$, $\overline{y}$, $\overline{xy}$, $\overline{x^2}$ en $\overline{y^2}$ nodig.
        Deze bepalen we aan de hand van de volgende rekentabel:
        \begin{center}
            \begin{tabular}{ccccc}
                \toprule
                    $x$ & $y$ & $xy$ & $x^2$ & $y^2$ \\
                \midrule
    		        $14.2$ & $43.5$ & $617.7$ & $201.64$ & $1892.25$ \\
                    $27.8$ & $26.2$ & $728.36$ & $772.84$ & $686.44$ \\
                    $33.5$ & $41.1$ & $1376.85$ & $1122.25$ & $1689.21$ \\
                    $46.1$ & $64.4$ & $2968.84$ & $2125.21$ & $4147.36$ \\
                    $52.4$ & $28.2$ & $1477.68$ & $2745.76$ & $795.24$ \\
                    $61.9$ & $45.6$ & $2822.64$ & $3831.61$ & $2079.36$ \\
                    $74.3$ & $66.8$ & $4963.24$ & $5520.49$ & $4462.24$ \\
                    $85.6$ & $62.3$ & $5332.88$ & $7327.36$ & $3881.29$ \\
                    $93.2$ & $89.5$ & $8341.4$ & $8686.24$ & $8010.25$ \\
                    $108.7$ & $71.2$ & $7739.44$ & $11815.69$ & $5069.44$ \\
                    $119.4$ & $94.6$ & $11295.24$ & $14256.36$ & $8949.16$ \\
                \midrule
                    $\overline{x} = 65.1909$ & $\overline{y} = 57.5818$ & $\overline{xy} = 4333.1155$ & $\overline{x^2} = 5309.5864$ & $\overline{y^2} = 3787.4764$ \rubric{4} \\ 
                \bottomrule
            \end{tabular}
        \end{center}
        \par\medskip
        De correlatiecoëfficiënt van Pearson is dus gelijk aan
        \begin{align*}
            r(x,y)  &= \frac{ \overline{x \cdot y} - \overline{x} \cdot \overline{y} }{ \sqrt{ (\overline{x^2} - \overline{x}^2) \cdot (\overline{y^2} - \overline{y}^2) } }\\
                    &= \frac{ 4333.1155 - 65.1909 \cdot 57.5818 }{ \sqrt{ (5309.5864 - 65.1909^2) \cdot (3787.4764 - 57.5818^2) } } \\
                    &= \frac{579.3044}{707.1016} \\
                    &\approx 0.8193.\rubric{2}
        \end{align*}
        De correlatiecoëfficiënt $r \approx 0.8193$ ligt ver van $0$ af en is positief.
        Dit houdt in dat er een sterke positieve (stijgende) trend is. \rubric{1}
        Een hoger munitieverbruik correleert sterk met een hoger percentage succesvol uitgeschakelde doelen. \rubric{1}

        (Maximaal $2$ punten indien de student een negatieve waarde heeft berekend voor $r(x, y)$ en inziet dat dit niet rijmt met het spreidingsdiagram.
        $0$ punten indien de student dit besef niet laat zien in zijn/haar uitwerking.)
    }

    \subquestion{6}{
        Bereken de coëfficiënten $a$ en $b$ van de regressielijn $\hat{Y} = a + b \cdot X$.
        Bepaal aan de hand van de regressielijn een statistisch verantwoorde voorspelling voor het 
        percentage succesvol geneutraliseerde doelen bij een operatie waarbij $80$ munitiekisten worden verschoten.
    }
    \solution{
        Om de coëfficiënten $a$ en $b$ van de regressielijn te berekenen, gaan we de rekentabel van vraag (a) hergebruiken.
        Er volgt namelijk dat
        \begin{align*}
            b &= \frac{\overline{xy} - \overline{x} \cdot \overline{y}}{\overline{x^2} - (\overline{x})^2} \\
              &= \frac{4333.1155 - 65.1909 \cdot 57.5818}{5309.5864 - (65.1909)^2} \\
              &= \frac{579.3044}{1059.7317} \approx 0.5467 \rubric{2} \\
            a &= \overline{y} - b \cdot \overline{x} \\
              &= 57.5818 - 0.5467 \cdot 65.1909 \\
              &\approx 21.9451. \rubric{2}
        \end{align*}
        De formule van de regressielijn behorende bij deze steekproef is dus gelijk aan $\hat{Y} = 21.9451 + 0.5467 \cdot X$. \rubric{1}
        Een statistisch verantwoorde voorspelling voor het percentage succesvol geneutraliseerde doelen bij een operatie met $80$ verbruikte munitiekisten
        vinden we door $X = 80$ in te vullen:
        Dit geeft een waarde van $\hat{Y}_0 = 21.9451 + 0.5467 \cdot 80 \approx 65.6772$, oftewel 
        onze voorspelling zegt dat bijna twee derde van de doelen succesvol worden geneutraliseerd. \rubric{1} 

        (Geen punten indien $b < 0$ zonder besef dat dit niet kan kloppen. \\
        Maximaal $2$ punten indien de coëfficiënten dusdanige waarden hebben dat deze niet logischerwijs een passende lijn vormen door de steekproefdata.)
    }

    \subquestion{8}{
        Bereken een \SI{95}{\percent}-betrouwbaarheidsinterval voor het verwachte succespercentage van geneutraliseerde doelen bij een operatie
        waarbij $80$ munitiekisten worden verschoten.
        Rond het interval af op gehele percentages zodat de betrouwbaarheid van het interval gewaarborgd blijft.
    }
    \solution{
        In opdracht (d) hebben we een puntschatting van $\hat{Y}_0 = 21.9451 + 0.5467 \cdot 80 \approx 65.6772$ bepaald.
        Daarnaast kunnen we de standaardafwijking $\sigma$ van de storingsterm $\varepsilon$ schatten:
        \begin{align*}
            s_{\varepsilon} &= \sqrt{ \frac{n}{n-2} \cdot \left( \overline{y^2} - a \cdot \overline{y} - b \cdot \overline{xy} \right) } \\ 
                               &= \sqrt{ \frac{11}{9} \cdot \left( 3787.4764 - 21.9451 \cdot 57.5818 - 0.5467 \cdot 4333.1155 \right) } \\ 
                               &\approx 13.7698.\rubric{2}
        \end{align*}

        Vervolgens kunnen we een puntschatting berekenen van de standaardafwijking van de verwachtingswaarde van $Y$ voor gegeven $X = x_0$:
        \begin{align*}
        s_{\mu}  &= s_{\varepsilon} \cdot \sqrt{ \frac{1}{n} \cdot \left( 1 + \frac{(x_0 - \overline{x})^2}{\overline{x^2} - \overline{x}^2} \right) } \\
                    &= 13.7698 \cdot \sqrt{ \frac{1}{11} \cdot \left( 1 + \frac{(80 - 65.1909)^2}{5309.5864 - 65.1909^2} \right) } \\
                    &\approx 4.5612.\rubric{2}
        \end{align*}

        Omdat we de standaardafwijkingen geschat hebben en de storingstermen normaal verdeeld zijn, moeten we werken met de $t$-verdeling met $\df = n - 2 = 9$ vrijheidsgraden.
        De $t$-waarde die hoort bij een betrouwbaarheidsniveau $\alpha = 0.05$ is gelijk aan
        \[
            t = \invt(\opp = 1 - \alpha / 2, \df = n - 2) = \invt(\opp = 0.975, \df = 9) \approx 2.2622.\rubric{1}
        \]
        Het \SI{95}{\percent}-betrouwbaarheidsinterval voor de gemiddelde $Y$ voor gegeven $X = 80$ kan dus worden beschreven door
        \begin{align*}
            &[y_0 - t \cdot s_{\mu}, y_0 - t \cdot s_{\mu}] \\
            &= [65.6772 - 2.2622 \cdot 4.5612, 65.6772 + 2.2622 \cdot 4.5612] \\
            &\approx [55.3592, 75.9953]. \rubric{2}
        \end{align*}
        
        Met \SI{95}{\percent} betrouwbaarheid ligt het gemiddelde van $Y$ voor gegeven $X = 80$ tussen $55.3592$ en $75.9953$ procent.       
        Om de betrouwbaarheid te waarborgen, moeten we het interval naar buiten afronden, dus $[55, 76]$. \rubric{1}
    }
\end{question}